\documentclass{article}%
\usepackage[T1]{fontenc}%
\usepackage[utf8]{inputenc}%
\usepackage{lmodern}%
\usepackage{textcomp}%
\usepackage{lastpage}%
\usepackage{geometry}%
\geometry{margin=2.2cm}%
%
\usepackage{amsmath}%
\usepackage{amssymb}%
\usepackage{booktabs}%
\usepackage{longtable}%
\usepackage{array}%
\usepackage{hyperref}%
\hypersetup{hidelinks}%
\title{Soluções\\[4pt]\large Comparação com o gabarito Gabarito oficial AFA 2027, Versão A\thanks{Gerado com gpt-resolve: \url{https://github.com/lgabs/gpt-resolve}}}%
\date{29/10/2026}%
%
\begin{document}%
\normalsize%
\maketitle%
\section*{Resumo}\textbf{Nota média: 10,0/10} (16 de 16 itens com a resolução correta; cada item vale o mesmo dentro da questão).\par\medskip%
\begin{longtable}{@{}c c p{0.78\textwidth}@{}}\toprule Questão & Nota & Resumo \\ \midrule
49 & 10,0 & A resolução e o gabarito batem: os dois marcam a alternativa A. Refiz a análise conferindo a figura original e cheguei ao mesmo resultado. De 0 a t, a partícula A tem v < 0 e a < 0, então o movimento é retrógrado e acelerado. B é progressivo e acelerado, o que derruba (b) e (c). Em C, a = 0, mas o sinal de v$_{0}$ é desconhecido, então (d) é falsa. A resolução está correta e bem justificada, e as dúvidas que sobram são pontos menores. \\
50 & 10,0 & A resolução está correta e coincide com o gabarito oficial (C, 2,4 $^\circ$C). O caminho físico está certo: lançamento a 45$^\circ$ com v0$^2$ = 3gh, conservação de energia incluindo o ganho de energia potencial no plano ($\frac12$kx$^2$ = 2,5mgh), Ec no topo = 0,60 J e metade dela aquecendo a esfera, $\Delta$T = 2,4 $^\circ$C. Conferi os números em python e não há divergência. Só ficam observações menores sobre a leitura da figura e a escala. \\
51 & 10,0 & A resolução está correta e coincide com o gabarito oficial (C), com 0,6 s. O máximo de Ec em L = 110 cm marca o equilíbrio, o que dá d = 0,10 m e T = 2$\pi$$\surd$(d/g) = 0,6 s. A verificação por energia (k = 100 N/m, m = 1,0 kg) confirma o resultado. Não há erro nem ambiguidade relevante. \\
52 & 10,0 & A resolução e o gabarito oficial (Versão A) concordam na alternativa D, e a física está correta. Sem escorregamento, v = $\omega$\_D R\_D = $\omega$\_T R\_T. Como R\_T = 2R\_D, $\omega$\_D = 2$\omega$\_T em todo instante, e os dois gráficos são trapézios com os mesmos instantes t1, t2 e t3, com a roda dianteira no patamar mais alto. Conferi as alternativas no PDF original, e só a (d) mostra essa configuração. Não há erro na resolução nem no gabarito. \\
53 & 10,0 & A resolução e o gabarito oficial (Versão A) concordam: alternativa D. A resposta está correta, porque a sustentação vem da diferença de pressão entre as faces da asa, causada pela diferença de velocidade do escoamento, com pressão menor onde a velocidade é maior. A física da resolução e a refutação das alternativas a, b e c estão certas, e as contas de ordem de grandeza conferem. O único problema é um erro de digitação no LaTeX (\textbackslash{}emph\{antes* sem fechar a chave), que não afeta o conteúdo. \\
54 & 10,0 & A resolução e o gabarito oficial (Versão A) concordam: a resposta da questão 54 é D. A física da resolução está certa. O pano condensa o vapor, a pressão sobre a água cai e a água, ainda quente, volta a ferver abaixo de 100 $^\circ$C até se resfriar o bastante. Os valores numéricos citados conferem com tabelas e com a equação de Antoine. Também está correto o descarte de a), b) e c). Não há erro relevante. O único ponto de redação é o termo 'pressão de vapor' na própria alternativa d). \\
55 & 10,0 & A resolução está correta e bate com o gabarito oficial: alternativa B, t' = t$\cdot$(1 + $\alpha$$\Delta$$\theta$)\^{}\{1/2\}. A aceleração g$\cdot$sen $\beta$ não muda com o aquecimento e, partindo do repouso, o tempo é proporcional a $\surd$L. Com L' = L$_{0}$(1 + $\alpha$$\Delta$$\theta$), sai o resultado. A transcrição confere com o original. Não há divergência a resolver. \\
56 & 10,0 & Questão 56: a resolução e o gabarito oficial (versão A) dão a mesma resposta, alternativa B (20 m/s), e ela está correta. Pelo filtro de velocidades, v = E/B = 4,0/0,2 = 20 m/s. Como o espelho plano é fixo, a imagem se aproxima dele com essa mesma velocidade. O raciocínio da resolução está completo e mostra por que 40 m/s, a velocidade relativa entre imagem e partícula, está errado. A única ressalva é de pouco peso: o enunciado não diz que v é perpendicular a B, mas a resolução assume isso de forma explícita. \\
57 & 10,0 & A resolução e o gabarito oficial (Versão A) concordam: alternativa B, 39 $^\circ$C. Como a potência do chuveiro não muda, a vazão multiplicada pela elevação de temperatura é constante. Então $\Delta$$\theta$$_{2}$ = 20$\cdot$182/130 = 28 $^\circ$C e a água sai a 11 + 28 = 39 $^\circ$C. A física e as contas estão corretas, e as únicas ressalvas são hipóteses implícitas razoáveis. \\
58 & 10,0 & A resolução marca (C), T = $\mu$(mg + BiL)/($\mu$ cos $\theta$ + sen $\theta$), igual ao gabarito oficial da AFA 2027, versão A. A física está correta: a força magnética aponta para baixo, o atrito em C aponta para a direita e, na iminência, f = $\mu$N. A resolução também confere o resultado pelos momentos, que dão $\mu$ = tg $\theta$ e T = W/(2 cos $\theta$), e isso bate com a expressão (c). Não há erro nem na resolução nem no gabarito. A única ressalva é que o enunciado tem mais condições do que incógnitas, o que não muda a resposta. \\
59 & 10,0 & A resolução está correta e bate com o gabarito oficial (C). Como os momentos são opostos e de mesmo módulo, os dois fótons têm a mesma energia. A frequência mínima corresponde ao par em repouso, e cada fóton leva m c$^2$ = 8,1$\cdot$10\^{}{-}14 J. Daí f = m c$^2$/h $\approx$ 1,2$\cdot$10\^{}20 Hz, de ordem 10\^{}20 Hz, o que conferi em python3. Restam só observações de redação, sem impacto na nota. \\
60 & 10,0 & A resolução e o gabarito coincidem: a alternativa correta é a (A), -4,0 $\mu$C, e a resolução está certa. A física está correta. Q\_C repele Q\_B, logo Q\_C < 0, e Q\_A atrai Q\_B, logo Q\_A = +1 $\mu$C. O equilíbrio horizontal F\_AB = F\_CB$\cdot$cos60$^\circ$ dá |Q\_C| = 2 $\mu$C, e a soma é -4 $\mu$C. Conferi as combinações de sinal numericamente. O único deslize é secundário, no comentário sobre o distrator (b). \\
61 & 10,0 & A resolução e o gabarito oficial da AFA 2027 (Versão A) marcam a mesma alternativa, a (A), e ela está correta. A dedução parte da relação relativística entre energia e momento, com E = m$_{0}$c$^2$ + K, e chega a p$^2$ = 2m$_{0}$K + K$^2$/c$^2$, portanto $\lambda$ = h/$\surd$(2m$_{0}$K + K$^2$/c$^2$). A análise de dimensões, os casos limite e a eliminação das outras alternativas estão certos. Não há nada a levar ao professor. \\
62 & 10,0 & A resolução e o gabarito oficial concordam na alternativa D, e as duas estão certas. O côncavo forma uma imagem real a 1,5R, com A1 = -2. Essa imagem fica a 0,5R do convexo e serve de objeto real para ele. O convexo forma uma imagem virtual R/4 atrás dele, com A2 = +1/2. O aumento total é -1, então a imagem final é virtual, invertida em relação ao objeto e com metade da altura da imagem do côncavo. A resolução está completa e examina cada alternativa corretamente. \\
63 & 10,0 & A resolução está correta e coincide com o gabarito oficial (B). Ela tira $\mu$ = tg$\theta$ da descida com velocidade constante e calcula v\_B$^2$ = 2gH(1 - tg$\theta$/tg$\alpha$). Na colisão elástica com a massa 2m, a partícula 2 sai com v$_{2}$' = 2v/3, e a energia se conserva na rampa lisa. Daí h = (4/9)(1 - tg$\theta$/tg$\alpha$)H. Conferi a conta numericamente. Não há erro nem no gabarito nem na resolução; só existem hipóteses implícitas razoáveis. \\
64 & 10,0 & A resolução e o gabarito oficial concordam na alternativa (B), $\varepsilon$ = $\pi$ B\_V L$^2$ f, e o resultado está correto. A rotação da pá no plano horizontal só gera f.e.m. entre M e N por causa da componente vertical do campo. A componente B\_H produz força perpendicular à pá. A f.e.m. vale B\_V $\omega$L$^2$/2, com $\omega$ = 2$\pi$f. A dedução está correta e foi conferida por dois caminhos (força de Lorentz e lei de Faraday), além de análise dimensional. Não há erro a apontar. \\
\bottomrule\end{longtable}%
\section*{Questão 49}%
\begin{longtable}{@{}p{0.06\textwidth} p{0.33\textwidth} p{0.33\textwidth} p{0.1\textwidth} p{0.08\textwidth}@{}}\toprule Item & Resolução & Gabarito & Veredito & Certo \\ \midrule
49 & (A) A é acelerado e retrógrado & A & igual & ambos \\
\bottomrule\end{longtable}%
\begin{itemize}%
\item%
\textbf{49)} Conferi a figura no PDF original e ela confere com a transcrição. Em A, a curva espaço $\times$ tempo é côncava para baixo (a < 0) e tem o máximo num instante negativo. Por isso, de 0 a t, o espaço cai de um valor positivo até zero, com inclinação cada vez maior: v < 0 (retrógrado) e |v| crescente, com v$\cdot$a > 0 (acelerado). Em B, v = k$\tau$ com k > 0, então o movimento é progressivo e acelerado; (b) e (c) são falsas. Em C, a = 0, mas o gráfico não informa o sinal de v$_{0}$ (o movimento pode até ser repouso), então (d) é falsa por causa do 'necessariamente'. A resolução e o gabarito dão a mesma resposta, e ela está correta.%
\end{itemize}%
\paragraph{Pontos de dúvida}%
\begin{itemize}%
\item%
\textit{(solution)} Em B, v = 0 no instante $\tau$ = 0. A resolução diz, com razão, que esse repouso momentâneo não muda a classificação (progressivo e acelerado) no intervalo, mas a rigor isso só vale para 0 < $\tau$ $\leq$ t.%
\item%
\textit{(solution)} A posição do vértice da parábola de A (em tempo negativo) foi lida no gráfico. Mesmo sem supor MUV, a concavidade para baixo (a < 0) e a curva decrescente de 0 a t (v < 0) já bastam para concluir que o movimento é acelerado e retrógrado.%
\end{itemize}%
\section*{Questão 50}%
\begin{longtable}{@{}p{0.06\textwidth} p{0.33\textwidth} p{0.33\textwidth} p{0.1\textwidth} p{0.08\textwidth}@{}}\toprule Item & Resolução & Gabarito & Veredito & Certo \\ \midrule
50 & (C) 2,4 $^\circ$C & C & igual & ambos \\
\bottomrule\end{longtable}%
\begin{itemize}%
\item%
\textbf{50)} Reli o enunciado e a figura no PDF original (p. 2). A linha tracejada horizontal passa pela posição inicial da esfera, e h é medido a partir dela até o ponto de lançamento. Refiz as contas em python. O lançamento a 45$^\circ$ com deslocamento horizontal até o topo d = v0$^2$/(2g) = 1,5h dá v0$^2$ = 3gh. Pela conservação da energia, $\frac12$kx$^2$ = 2,0 J = mgh + 1,5mgh, logo mgh = 0,80 J (h = 80 m e v0$^2$ = 2400 m$^2$/s$^2$). No topo sobra só a componente horizontal da velocidade: Ec = $\frac12$m(v0$^2$/2) = 0,60 J. Metade disso, 0,30 J, é igual a m$\cdot$c$\cdot$$\Delta$T, o que dá $\Delta$T = 0,30/(10$^{-}$$^3$$\cdot$125) = 2,4 $^\circ$C. A resolução e o gabarito batem e estão corretos.%
\end{itemize}%
\paragraph{Pontos de dúvida}%
\begin{itemize}%
\item%
\textit{(statement)} A altura h vem da leitura da figura: ela é medida da linha tracejada, que passa pela posição inicial da esfera (mola comprimida). Se h fosse medido a partir da posição de mola relaxada, seria preciso somar x$\cdot$sen45$^\circ$ $\approx$ 0,14 m de subida. Com h = 80 m, isso muda $\Delta$T só na terceira casa decimal ($\approx$ 2,398 $^\circ$C), sem alterar a alternativa.%
\item%
\textit{(statement)} Os dados levam a h = 80 m, muito maior que os 20 cm de compressão da mola. A figura, portanto, está fora de escala. Isso não afeta o resultado, mas pode chamar a atenção de quem confere.%
\item%
\textit{(solution)} A explicação de onde vêm os distratores (9,6; 4,8; 1,2) é especulativa. Ela não interfere no raciocínio principal, que está correto e completo, com verificação da energia total no topo (0,60 + 0,80 + 0,60 = 2,0 J).%
\end{itemize}%
\section*{Questão 51}%
\begin{longtable}{@{}p{0.06\textwidth} p{0.33\textwidth} p{0.33\textwidth} p{0.1\textwidth} p{0.08\textwidth}@{}}\toprule Item & Resolução & Gabarito & Veredito & Certo \\ \midrule
51 & (C) 0,6 s & C & igual & ambos \\
\bottomrule\end{longtable}%
\begin{itemize}%
\item%
\textbf{51)} A transcrição confere com o original (p. 19 da Versão A): mola de comprimento natural 100 cm e gráfico de Ec $\times$ L com zeros em 100 e 120 cm e máximo de 0,5 J em 110 cm. Refiz a conta. Ec é máxima onde a resultante é nula, logo o equilíbrio está em L = 110 cm, e a deformação estática é d = 0,10 m. Daí k/m = g/d = 100 s$^{-}$$^2$ e T = 2$\pi$$\surd$(d/g) = 2$\cdot$3,0$\cdot$0,10 = 0,6 s. Conferi também pela energia: a amplitude é A = 0,10 m e Ec,máx = kA$^2$/2 = 0,5 J, o que dá k = 100 N/m. Então m = kd/g = 1,0 kg e T = 2$\cdot$3$\cdot$$\surd$(1/100) = 0,6 s. Pela conservação de 100 a 120 cm, mg$\cdot$0,20 = 2,0 J = k(0,20)$^2$/2. Todas as contas da resolução estão certas, e a resposta coincide com o gabarito (C).%
\end{itemize}%
\paragraph{Pontos de dúvida}%
\begin{itemize}%
\item%
\textit{(solution)} Os dados 110 cm e 0,5 J foram lidos no gráfico. Como estão marcados com linhas tracejadas e rótulos explícitos, não há ambiguidade real. Para o período basta a abscissa do máximo, então o valor de 0,5 J serve só para a verificação de consistência.%
\item%
\textit{(solution)} A resolução usa $\pi$ = 3,0, como manda a prova, e chega a 0,6 s. Ela observa corretamente que com $\pi$ exato T $\approx$ 0,63 s, valor que continua mais próximo de (c).%
\end{itemize}%
\section*{Questão 52}%
\begin{longtable}{@{}p{0.06\textwidth} p{0.33\textwidth} p{0.33\textwidth} p{0.1\textwidth} p{0.08\textwidth}@{}}\toprule Item & Resolução & Gabarito & Veredito & Certo \\ \midrule
52 & (D): $\omega$\_D = 2$\omega$\_T em todo instante; dois trapézios com os mesmos instantes t1, t2, t3, e $\omega$\_D no patamar 2$\omega$, $\omega$\_T no patamar $\omega$ & D & igual & ambos \\
\bottomrule\end{longtable}%
\begin{itemize}%
\item%
\textbf{52)} A resolução bate com o gabarito. Conferi a física: sem escorregamento vale v = $\omega$R para cada roda, e as duas rodas têm a velocidade do chassi v(t). Com R\_T = 2R\_D, $\omega$\_T = v/(2R\_D) = $\omega$\_D/2. Então $\omega$\_D = 2$\omega$\_T e as duas são proporcionais a v(t), o que dá trapézios com os mesmos instantes, e a roda dianteira (a menor) fica no patamar mais alto. Pela página 3 do PDF, a alternativa (d) é a única que tem $\omega$\_D em 2$\omega$ e $\omega$\_T em $\omega$. A (a) inverte os rótulos, a (b) põe $\omega$\_D = $\omega$\_T e a (c) também deixa $\omega$\_T acima de $\omega$\_D. A transcrição do enunciado e das alternativas confere com o original.%
\end{itemize}%
\paragraph{Pontos de dúvida}%
\begin{itemize}%
\item%
\textit{(statement)} Nos gráficos (a) e (d), as alternativas se diferenciam só pela posição dos rótulos $\omega$\_D/$\omega$\_T, que são ligados às curvas por setas curvas na parte de descida, perto de t3. Na imagem da página 3, (d) tem $\omega$\_D no trapézio de cima (2$\omega$) e $\omega$\_T no de baixo ($\omega$), e (a) tem o contrário. A leitura é clara, mas depende de seguir as setas na figura. Não consegui ampliar o PDF em alta resolução porque o ambiente bloqueou a renderização.%
\item%
\textit{(solution)} A resolução supõe, sem dizer, que as duas rodas rolam sem escorregar ao mesmo tempo e que os eixos têm a mesma velocidade, por estarem no mesmo chassi rígido em movimento retilíneo. A hipótese é razoável e é a que o enunciado pretende.%
\end{itemize}%
\section*{Questão 53}%
\begin{longtable}{@{}p{0.06\textwidth} p{0.33\textwidth} p{0.33\textwidth} p{0.1\textwidth} p{0.08\textwidth}@{}}\toprule Item & Resolução & Gabarito & Veredito & Certo \\ \midrule
Questão 53 (única) & D & D & igual & ambos \\
\bottomrule\end{longtable}%
\begin{itemize}%
\item%
\textbf{Questão 53 (única))} Conferi o enunciado no PDF (pág. 3, Versão A) e a transcrição está fiel. A alternativa (d) é a única que aplica Bernoulli corretamente: a sustentação é a resultante da diferença de pressão entre as faces da asa, e a pressão é menor onde a velocidade é maior. As outras alternativas estão erradas: (a) é a falácia do tempo de trânsito igual; (b) atribui a queda de pressão a uma 'conversão direta' de energia cinética em potencial; (c) descreve só o empuxo hidrostático, que é desprezível e existe mesmo com o avião parado. Refiz as estimativas numéricas da resolução e todas conferem: $\rho$g$\Delta$y $\approx$ 2,4 Pa; $\frac12$$\rho$v$^2$ $\approx$ 2940 Pa; uma diferença de $\pm$5\% nas velocidades dá $\Delta$p $\approx$ 588 Pa; carga alar 1000$\cdot$10/16 $\approx$ 625 Pa.%
\end{itemize}%
\paragraph{Pontos de dúvida}%
\begin{itemize}%
\item%
\textit{(solution)} Erro de LaTeX na análise da alternativa (a): está escrito '\textbackslash{}emph\{antes* ao bordo de fuga', com '*' no lugar de '\}'. A chave fica aberta, o que pode quebrar a compilação ou deixar em ênfase todo o texto seguinte. Não muda o conteúdo físico.%
\item%
\textit{(solution)} A resolução usa a mesma constante de Bernoulli para linhas de corrente diferentes (por cima e por baixo da asa). A justificativa dada é que o escoamento ao longe é uniforme. Isso vale em escoamento estacionário e invíscido com energia total uniforme a montante, e é adequado ao nível da prova. Os valores de $\rho$, espessura da asa, velocidade e carga alar são estimativas de ordem de grandeza que o enunciado não fornece. A resposta não depende deles.%
\end{itemize}%
\section*{Questão 54}%
\begin{longtable}{@{}p{0.06\textwidth} p{0.33\textwidth} p{0.33\textwidth} p{0.1\textwidth} p{0.08\textwidth}@{}}\toprule Item & Resolução & Gabarito & Veredito & Certo \\ \midrule
54 & (D): o vapor condensa por causa da troca de calor com o pano molhado, a pressão sobre o líquido cai e a água, ainda quente, volta a ferver abaixo de 100 $^\circ$C. & D & igual & ambos \\
\bottomrule\end{longtable}%
\begin{itemize}%
\item%
\textbf{54)} A transcrição confere com o original (p. 21 do caderno, Versão A). A física é a da clássica 'fervura com água fria' (experimento de Franklin). A temperatura de ebulição é aquela em que p\_sat(T) iguala a pressão sobre o líquido. O pano resfria a parede em contato com o vapor, o vapor condensa e a pressão no balão cai para perto de p\_sat da parede fria. Como a água esfria devagar, ela fica com p\_sat(T\_água) maior que a pressão interna e ferve. A própria vaporização a resfria até a fervura parar, o que explica o 'até certa temperatura'. Conferi os valores de tabela citados com a equação de Antoine: p\_sat(25 $^\circ$C) $\approx$ 3,2 kPa, p\_sat(80 $^\circ$C) $\approx$ 47 kPa, p\_sat(90 $^\circ$C) $\approx$ 70 kPa e p\_sat(100 $^\circ$C) $\approx$ 101 kPa. A estimativa $\rho$gh $\approx$ 1 kPa também está correta. As razões dadas para descartar a), b) e c) estão certas: não há fonte de calor, a água só perde energia, e condensar não faz o líquido ferver.%
\end{itemize}%
\paragraph{Pontos de dúvida}%
\begin{itemize}%
\item%
\textit{(statement)} A alternativa d) fala em 'diminuindo a pressão de vapor'. Rigorosamente, o que cai é a pressão do gás (vapor) sobre o líquido. A pressão de vapor da água só depende da temperatura dela e continua alta. É um uso informal do termo, que não deixa dúvida sobre a alternativa correta, mas vale comentar com os alunos.%
\item%
\textit{(solution)} A resolução supõe que a pressão interna tende a p\_sat na temperatura da parede fria (cerca de 3 kPa) e que o ar residual é desprezível. Na prática a pressão fica acima disso, pelo ar que sobra e pela evaporação contínua da água quente. A conclusão não muda, porque basta que a pressão interna fique abaixo de p\_sat(T\_água).%
\end{itemize}%
\section*{Questão 55}%
\begin{longtable}{@{}p{0.06\textwidth} p{0.33\textwidth} p{0.33\textwidth} p{0.1\textwidth} p{0.08\textwidth}@{}}\toprule Item & Resolução & Gabarito & Veredito & Certo \\ \midrule
55 & (B) t$\cdot$(1 + $\alpha$$\cdot$$\Delta$$\theta$)\^{}\{1/2\} & B & igual & ambos \\
\bottomrule\end{longtable}%
\begin{itemize}%
\item%
\textbf{55)} Conferi com a página 4 do PDF: a transcrição está fiel (calha lisa, homogênea, mesma inclinação depois do aquecimento, mesmas alternativas). Refiz a física. Numa calha lisa a aceleração é a = g$\cdot$sen $\beta$, que não muda com o aquecimento. Partindo do repouso, L = a$\cdot$t$^2$/2, então t é proporcional a $\surd$L. Depois do aquecimento, L' = L$_{0}$(1 + $\alpha$$\Delta$$\theta$), logo t' = t$\cdot$(1 + $\alpha$$\Delta$$\theta$)\^{}\{1/2\}. A resolução e o gabarito dão a mesma resposta, e ela está correta. A análise das alternativas também está certa: (a) leva a t' $\to$ 0 quando $\Delta$$\theta$ $\to$ 0; (c) valeria se t fosse proporcional a L; (d) se t fosse proporcional a L$^2$.%
\end{itemize}%
\paragraph{Pontos de dúvida}%
\begin{itemize}%
\item%
\textit{(solution)} A resolução trata a partícula como ponto material que desliza. Na figura ela aparece como uma pequena esfera. Se ela rolasse, a aceleração seria uma fração constante de g$\cdot$sen $\beta$ (por exemplo, (5/7)g$\cdot$sen $\beta$ para uma esfera maciça). A razão t'/t seria a mesma, então a resposta não muda.%
\item%
\textit{(solution)} A resolução supõe, sem dizer, que a partícula percorre a calha inteira, do topo à base, e que a dilatação linear vale em todo o comprimento (calha homogênea). As duas hipóteses são razoáveis e vêm do próprio enunciado.%
\end{itemize}%
\section*{Questão 56}%
\begin{longtable}{@{}p{0.06\textwidth} p{0.33\textwidth} p{0.33\textwidth} p{0.1\textwidth} p{0.08\textwidth}@{}}\toprule Item & Resolução & Gabarito & Veredito & Certo \\ \midrule
56 & (B) 20 m/s & B & igual & ambos \\
\bottomrule\end{longtable}%
\begin{itemize}%
\item%
\textbf{56)} A transcrição confere com o original (p. 4 do PDF, versão A): E = 4,0 N/C, B = 2,0$\cdot$10$^{-}$$^{1}$ T, alternativas 10/20/40/80. Se a velocidade é constante, a força resultante é nula: qE = qvB, então v = E/B = 4,0/0,2 = 20 m/s. A trajetória é perpendicular ao espelho, que está fixo. Logo, a imagem se aproxima do espelho com a mesma velocidade da partícula, 20 m/s. O valor 40 m/s seria a velocidade relativa entre imagem e partícula, e a resolução descarta esse distrator corretamente. Resolução e gabarito estão de acordo e corretos.%
\end{itemize}%
\paragraph{Pontos de dúvida}%
\begin{itemize}%
\item%
\textit{(statement)} O enunciado não diz como v, E e B estão orientados entre si. Pela condição de força nula, E fica obrigatoriamente perpendicular a v, já que v$\times$B é perpendicular a v. Mas v poderia ter uma componente paralela a B, e aí só a componente perpendicular a B valeria E/B. A resposta 20 m/s pressupõe a montagem clássica do seletor de velocidades, com v perpendicular a E e a B. A resolução declara essa hipótese de forma explícita, e ela é a leitura pretendida pela banca.%
\end{itemize}%
\section*{Questão 57}%
\begin{longtable}{@{}p{0.06\textwidth} p{0.33\textwidth} p{0.33\textwidth} p{0.1\textwidth} p{0.08\textwidth}@{}}\toprule Item & Resolução & Gabarito & Veredito & Certo \\ \midrule
57 & (B) 39 $^\circ$C & B & igual & ambos \\
\bottomrule\end{longtable}%
\begin{itemize}%
\item%
\textbf{57)} O enunciado na página 5 do PDF (Versão A, pág. 22) confere com a transcrição. A potência P = U$^2$/R não muda, porque a parte elétrica não foi alterada e a resistividade não depende da temperatura. Por isso P = $\rho$c?$\Delta$$\theta$ e o produto ?$\cdot$$\Delta$$\theta$ fica constante. Daí 182$\cdot$20 = 130$\cdot$$\Delta$$\theta$$_{2}$, o que dá $\Delta$$\theta$$_{2}$ = 28 $^\circ$C e $\theta$$_{2}$ = 11 + 28 = 39 $^\circ$C. Conferi a conta em python: 28,0 e 39,0. A resolução e o gabarito concordam e estão corretos.%
\end{itemize}%
\paragraph{Pontos de dúvida}%
\begin{itemize}%
\item%
\textit{(solution)} A resolução supõe sem dizer que a densidade da água é constante e que a fração da potência entregue à água é a mesma nas duas vazões, ou seja, que as perdas não mudam. As duas hipóteses são razoáveis e o enunciado as sugere ao mandar desprezar a influência da temperatura.%
\item%
\textit{(solution)} Na conferência de ordem de grandeza, a resolução estima P $\approx$ 15 kW (4,18$\cdot$182$\cdot$20 $\approx$ 15,2 kW). É bem mais do que um chuveiro comum, que costuma ter de 5,5 a 7,5 kW. Isso indica que os dados do enunciado são irreais, mas a resposta não muda, e a própria resolução diz que a estimativa é só uma conferência.%
\end{itemize}%
\section*{Questão 58}%
\begin{longtable}{@{}p{0.06\textwidth} p{0.33\textwidth} p{0.33\textwidth} p{0.1\textwidth} p{0.08\textwidth}@{}}\toprule Item & Resolução & Gabarito & Veredito & Certo \\ \midrule
58 & (C) T = $\mu$(mg + BiL)/($\mu$ cos $\theta$ + sen $\theta$) & C & igual & ambos \\
\bottomrule\end{longtable}%
\begin{itemize}%
\item%
\textbf{58)} Conferi a figura no PDF, na página 5: o fio sai do teto à esquerda de A, então puxa A para cima e para a esquerda, e o atrito em C aponta para a direita. A força magnética é i(-x?)$\times$(-B$\hat z$) = -BiL $\hat y$, ou seja, aponta para baixo, e por isso a carga vertical total é W = mg + BiL. As condições f = T sen $\theta$, T cos $\theta$ + N = W e f = $\mu$N dão T = $\mu$W/($\mu$ cos $\theta$ + sen $\theta$), que é a alternativa (c). Pelo equilíbrio de momentos em torno de A, N = W/2, o que força $\mu$ = tg $\theta$ e T = W/(2 cos $\theta$). Testei numericamente com $\theta$ = 0,6 rad e $\mu$ = tg $\theta$: a expressão (c) dá o mesmo valor de W/(2 cos $\theta$), e a (b) dá outro valor. A resolução está certa e bate com o gabarito oficial.%
\end{itemize}%
\paragraph{Pontos de dúvida}%
\begin{itemize}%
\item%
\textit{(statement)} O sistema tem mais equações do que incógnitas. O equilíbrio de momentos dá N = W/2 e T cos $\theta$ = W/2, e então a iminência de deslizar só acontece se $\mu$ = tg $\theta$. A tração pode ser escrita de várias formas equivalentes: W/(2 cos $\theta$), $\mu$W/(2 sen $\theta$) ou $\mu$W/($\mu$ cos $\theta$ + sen $\theta$). O gabarito escolheu a forma que não usa os momentos. A resolução comenta isso no item 5, de forma correta.%
\item%
\textit{(statement)} A forma $\mu$W/(2 sen $\theta$) também é válida. Por isso a alternativa (a) só está errada por causa do sinal: ela traz (mg - BiL). Ela é o distrator de quem erra o sentido da força magnética. A resposta depende de ler corretamente que a corrente vai de C para A e que o campo entra na folha, o que faz a força apontar para baixo.%
\end{itemize}%
\section*{Questão 59}%
\begin{longtable}{@{}p{0.06\textwidth} p{0.33\textwidth} p{0.33\textwidth} p{0.1\textwidth} p{0.08\textwidth}@{}}\toprule Item & Resolução & Gabarito & Veredito & Certo \\ \midrule
59 & (C) 10\^{}20 Hz (f\_min = m c$^2$/h $\approx$ 1,2$\cdot$10\^{}20 Hz) & C & igual & ambos \\
\bottomrule\end{longtable}%
\begin{itemize}%
\item%
\textbf{59)} Conferi o enunciado na página 5 do PDF e ele bate com a transcrição. Como os dois fótons saem com momentos opostos, de mesmo módulo, eles também têm a mesma energia. A frequência é mínima quando o par está em repouso (K = 0). Nesse caso, cada fóton fica com m c$^2$ = 9,0$\cdot$10\^{}{-}31 $\cdot$ (3,0$\cdot$10\^{}8)$^2$ = 8,1$\cdot$10\^{}{-}14 J, que dá cerca de 0,51 MeV. Então f = 8,1$\cdot$10\^{}{-}14 / 6,6$\cdot$10\^{}{-}34 $\approx$ 1,23$\cdot$10\^{}20 Hz, com log10 $\approx$ 20,09, e a ordem de grandeza é 10\^{}20 Hz. Refiz a conta em python3. A resolução e o gabarito marcam a mesma alternativa, a C.%
\end{itemize}%
\paragraph{Pontos de dúvida}%
\begin{itemize}%
\item%
\textit{(solution)} Na análise das alternativas, a resolução diz que a (d) 'contraria a hipótese de frequência mínima'. Fisicamente, se o par tiver energia cinética muito grande, os fótons podem ter frequências maiores. A (d) só é descartada porque a pergunta é sobre a frequência mínima. O argumento continua certo, apenas a redação podia ser mais precisa.%
\item%
\textit{(statement)} O enunciado não diz que o par está em repouso, só que 'suas massas de repouso são convertidas em energia'. A frequência mínima vem da hipótese K = 0, que a resolução explicita. Mesmo um erro comum, como dar 2mc$^2$ a um único fóton, levaria a $\approx$ 2,5$\cdot$10\^{}20 Hz, ainda de ordem 10\^{}20. Por isso essa ambiguidade não muda a alternativa.%
\end{itemize}%
\section*{Questão 60}%
\begin{longtable}{@{}p{0.06\textwidth} p{0.33\textwidth} p{0.33\textwidth} p{0.1\textwidth} p{0.08\textwidth}@{}}\toprule Item & Resolução & Gabarito & Veredito & Certo \\ \midrule
60 & (A) -4,0 $\mu$C  (Q\_A = +1 $\mu$C, Q\_C = -2 $\mu$C) & A & igual & ambos \\
\bottomrule\end{longtable}%
\begin{itemize}%
\item%
\textbf{60)} Conferi a figura original (pág. 23 do caderno, Versão A). A força F sobre Q\_B aponta verticalmente para baixo, perpendicular a AB. A força de A sobre B é horizontal, então só F\_CB tem componente vertical. Para essa componente apontar para baixo, C tem de repelir B, logo Q\_C < 0. A componente horizontal de F\_CB vale F\_CB$\cdot$cos60$^\circ$ e aponta para +x. Para anulá{-}la, F\_AB tem de apontar de B para A, que é atração, logo Q\_A = +1 $\mu$C. O módulo sai de |Q\_A| = |Q\_C|/2, logo Q\_C = -2 $\mu$C. Soma: 1 - 3 - 2 = -4 $\mu$C. Testei em Python as 8 combinações (Q\_A = $\pm$1, Q\_C = $\pm$2, $\pm$0,5). Só (+1, -2) dá resultante puramente vertical para baixo, com F = (0; -5,196) em unidades de k/L$^2$.%
\end{itemize}%
\paragraph{Pontos de dúvida}%
\begin{itemize}%
\item%
\textit{(solution)} O comentário sobre os distratores tem uma imprecisão. A resolução diz que a alternativa b (-1,5) aparece ao inverter a projeção (|Q\_C| = |Q\_A|/2). Mas, com o sinal correto Q\_C = -0,5 $\mu$C, a soma seria 1 - 3 - 0,5 = -2,5. Para dar -1,5 é preciso também errar o sinal (Q\_C = +0,5 $\mu$C). Isso não afeta a resposta final.%
\end{itemize}%
\section*{Questão 61}%
\begin{longtable}{@{}p{0.06\textwidth} p{0.33\textwidth} p{0.33\textwidth} p{0.1\textwidth} p{0.08\textwidth}@{}}\toprule Item & Resolução & Gabarito & Veredito & Certo \\ \midrule
61 & (A) $\lambda$ = h / $\surd$(2m$_{0}$K + K$^2$/c$^2$) & A & igual & ambos \\
\bottomrule\end{longtable}%
\begin{itemize}%
\item%
\textbf{61)} Conferi o enunciado na página 6 do PDF e ele bate com a transcrição, inclusive as alternativas. Refiz a conta: E = m$_{0}$c$^2$ + K e E$^2$ = p$^2$c$^2$ + m$_{0}$$^2$c$^{4}$. Daí p$^2$c$^2$ = 2m$_{0}$c$^2$K + K$^2$, ou seja, p$^2$ = 2m$_{0}$K + K$^2$/c$^2$, e $\lambda$ = h/p dá a alternativa (a). A resolução também descarta as outras corretamente. Em (c) e (d), K/c tem dimensão de momento e 2m$_{0}$K de momento ao quadrado, então a soma é inconsistente. Em (b), o radicando fica negativo quando K < 2m$_{0}$c$^2$. Os limites clássico (h/$\surd$(2m$_{0}$K)) e ultrarrelativístico (hc/K) estão certos.%
\end{itemize}%
\section*{Questão 62}%
\begin{longtable}{@{}p{0.06\textwidth} p{0.33\textwidth} p{0.33\textwidth} p{0.1\textwidth} p{0.08\textwidth}@{}}\toprule Item & Resolução & Gabarito & Veredito & Certo \\ \midrule
62 & (D) Imagem virtual, invertida em relação ao objeto (A = -1) e com metade da altura da imagem do côncavo (A2 = +1/2) & D & igual & ambos \\
\bottomrule\end{longtable}%
\begin{itemize}%
\item%
\textbf{62)} Conferi o enunciado na p. 6 do PDF (p. 23 da prova), e ele coincide com a transcrição. Refiz as contas.\newline%
\newline%
No espelho côncavo: p1 = 0,75R e f1 = R/2, o que dá p1' = 1,5R (imagem real) e A1 = -2.\newline%
\newline%
I1 fica entre os espelhos, a 2R - 1,5R = 0,5R do convexo. Por isso funciona como objeto real para ele.\newline%
\newline%
No espelho convexo: f2 = -R/2, o que dá p2' = -R/4 (imagem virtual) e A2 = +1/2.\newline%
\newline%
O aumento total é A = A1$\cdot$A2 = -1. Portanto a imagem final é virtual, invertida em relação ao objeto e tem metade da altura de I1, que é a alternativa D.\newline%
\newline%
A alternativa (a) falha só porque a altura final é igual à do objeto, e não metade dela. As alternativas (b) e (c) dizem que a imagem é real, o que está errado.\newline%
\newline%
Confirmei as contas em python com frações exatas.%
\end{itemize}%
\paragraph{Pontos de dúvida}%
\begin{itemize}%
\item%
\textit{(statement)} O objeto está entre os espelhos, então também manda luz direto para o convexo, e o próprio objeto extenso pode bloquear parte dos raios que convergem do côncavo. O enunciado deixa claro que interessa só o caminho objeto $\to$ côncavo $\to$ convexo, e a resolução adota isso. Não muda a resposta.%
\item%
\textit{(solution)} A resolução usa a convenção de Gauss: f > 0 no côncavo, f < 0 no convexo e p' < 0 para imagem virtual. Ela declara essa convenção e a aplica de forma coerente. Também está certa ao tratar I1 como objeto real para o convexo, porque 1,5R < 2R e os raios divergem a partir de I1 antes de chegar ao convexo.%
\end{itemize}%
\section*{Questão 63}%
\begin{longtable}{@{}p{0.06\textwidth} p{0.33\textwidth} p{0.33\textwidth} p{0.1\textwidth} p{0.08\textwidth}@{}}\toprule Item & Resolução & Gabarito & Veredito & Certo \\ \midrule
63 & (B) h = (4/9)(1 - tg$\theta$/tg$\alpha$)$\cdot$H & B & igual & ambos \\
\bottomrule\end{longtable}%
\begin{itemize}%
\item%
\textbf{63)} Comparei a transcrição com a página 7 do PDF e ela confere com o original, inclusive as alternativas e a figura. Refiz a conta. Com velocidade constante na inclinação $\theta$, $\mu$ = tg$\theta$. Na inclinação $\alpha$, a = g(sen$\alpha$ - tg$\theta$$\cdot$cos$\alpha$) e L = H/sen$\alpha$, o que dá v\_B$^2$ = 2gH(1 - tg$\theta$/tg$\alpha$). Na colisão elástica de m com 2m em repouso, v$_{2}$' = 2v/3, e a energia cinética se conserva. Na rampa lisa, h = v$_{2}$'$^2$/(2g) = (4/9)(1 - tg$\theta$/tg$\alpha$)H. Conferi numericamente com $\alpha$ = 50$^\circ$, $\theta$ = 20$^\circ$ e H = 3,7: as duas expressões dão h $\approx$ 1,1422. A resposta é a alternativa B, a mesma da resolução e do gabarito.%
\end{itemize}%
\paragraph{Pontos de dúvida}%
\begin{itemize}%
\item%
\textit{(solution)} A resolução supõe que a partícula não perde energia na passagem pelo vértice B (mudança brusca de direção). O enunciado não diz isso de forma explícita, mas é a hipótese que o problema pede, e a resolução a declara.%
\item%
\textit{(statement)} A resolução usa o mesmo $\mu$ (cinético) com a inclinação $\alpha$. Isso se apoia na expressão 'sobre o mesmo plano' e é uma hipótese razoável. O enunciado também não diz que a colisão é frontal (unidimensional), mas isso está implícito pela geometria.%
\end{itemize}%
\section*{Questão 64}%
\begin{longtable}{@{}p{0.06\textwidth} p{0.33\textwidth} p{0.33\textwidth} p{0.1\textwidth} p{0.08\textwidth}@{}}\toprule Item & Resolução & Gabarito & Veredito & Certo \\ \midrule
64 & (B) $\pi$ B\_V L$^2$ f & B & igual & ambos \\
\bottomrule\end{longtable}%
\begin{itemize}%
\item%
\textbf{64)} Conferi o enunciado na página 7 do PDF: a transcrição está fiel, tanto o texto quanto a figura e as alternativas. A física está certa. A pá gira num plano horizontal, então v e d? são horizontais. Com B\_H, o produto v $\times$ B\_H fica vertical, perpendicular à pá, e não contribui para a f.e.m. Com B\_V, v $\times$ B\_V fica radial, de módulo $\omega$rB\_V. Integrando ao longo da pá: $\varepsilon$ = $\int$$_{0}$? $\omega$rB\_V dr = B\_V $\omega$L$^2$/2. Com $\omega$ = 2$\pi$f, isso dá $\pi$ B\_V L$^2$ f, que é a alternativa (b). A conferência por Faraday também fecha: o fluxo de B\_V através do disco $\pi$L$^2$, varrido em T = 1/f, dá o mesmo valor. A estimativa numérica também confere: 3$\cdot$5$\cdot$10$^{-}$$^{5}$$\cdot$25$\cdot$5 $\approx$ 1,9$\cdot$10$^{-}$$^2$ V.%
\end{itemize}%
\paragraph{Pontos de dúvida}%
\begin{itemize}%
\item%
\textit{(solution)} O argumento usado para descartar a alternativa (c) diz que ela 'só valeria se o campo fosse inteiramente vertical'. Está correto, mas é redundante: nesse caso B\_H = 0 e (c) coincide com (b). A redação não afeta a resposta.%
\end{itemize}%
\end{document}